\documentclass[../../main.tex]{subfiles}
\begin{document}

\begin{center}
  {\fontsize{48pt}{58pt}\selectfont \textbf{2014}\normalsize\textbf{年度}}
\end{center}

\begin{table}[htb]
  \centering
  \caption{出題テーマと難易度}
  \begin{tabular}{|c|p{9cm}|c|}
    \hline
    \textbf{問題番号} & \textbf{テーマ} & \textbf{難易度} \\
    \hline
    第1問 & 空間の直線への距離の二乗和 & \\
    \hline
    第2問 & 三角形上を動く2粒子の一致確率 & \\
    \hline
    第3問 & 面積最大の三角形と角 & \\
    \hline
    第4問 & 有理関数の値域と係数の領域 & \\
    \hline
    第5問 & 立方数の和の整除性と最小値 & \\
    \hline
    第6問 & 円と双曲線で囲まれた面積 & \\
    \hline
  \end{tabular}
\end{table}

\begin{table}[htb]
  \centering
  \caption{解答}
  \renewcommand{\arraystretch}{1.4}
  \begin{tabular}{|c|c|p{9cm}|}
    \hline
    \textbf{問題番号} & \textbf{小問} & \textbf{答え} \\
    \hline
    第1問 & -- & $P=(1,0,-2)$で最小値$7$ \\
    \hline
    第2問 & -- & $p(n)=\dfrac{1}{3}\left(1+\dfrac{1}{2}4^{1-n}\right)$ \\
    \hline
    第3問 & -- & $\cos\angle B=\dfrac{\sqrt{33}-1}{8}$ \\
    \hline
    第4問 & -- & $\dfrac{a^2+2a-3}{4}\le b\le\dfrac{-a^2+2a+3}{4}$（境界を含む） \\
    \hline
    第5問 & -- & $(a,b)=(13,14),(14,13)$で最小値$365$ \\
    \hline
    第6問 & -- & $\dfrac{2\pi}{3}-\log(2+\sqrt{3})$ \\
    \hline
  \end{tabular}
\end{table}

\paragraph{試験形態}
6問．試験時間は要確認．

\end{document}
